\section{社区感染上限与医疗需求的联系}
\label{sec:sup:medical-link}
本节在正文传播模型之外补充一个医疗需求关系，说明如何在附加假设下根据容量选取社区感染上限。

两类感染仓室的新增流入之和为
\[
\frac{\beta c(t)(1-q(t))S(t)I(t)}{N}
+\frac{\beta c(t)q(t)S(t)I(t)}{N}
=\frac{\beta c(t)S(t)I(t)}{N}.
\]
若全过程满足 $0\le c(t)\le c_0$ 和 $I(t)\le\eta$，由无回流模型中的 $S(t)\le S_0$，得到
\[
0\le\frac{\beta c(t)S(t)I(t)}{N}
\le\frac{\beta c_0S_0}{N}\eta.
\]
这一上界包含进入 $I$ 和 $I_q$ 的全部新增感染。

设非负可测函数 $h(a)$ 表示一例新感染者在感染后第 $a$ 天所需的平均 ICU 床位数，并假定它对两类病例、不同感染时点及所比较的干预分配共同适用。令
\[
\int_0^\infty h(a)\,da=pL<\infty,
\]
其中 $p\in(0,1]$ 为需要 ICU 的概率，$L>0$ 为需要 ICU 者累计占用时长的均值；$h$ 可包含感染至入院的延迟。根据既往流入与停留过程计算床位占用的建模方法可参见文献\cite{Baas2021Occupancy}。

初始时刻之前已经感染者此后的平均需求记为 $D_0(t)$，并假定各策略均满足 $0\le D_0(t)\le D_{\mathrm{init}}<\infty$。该项包括初始病例以后才发生的入院和后续占用，也包括已从传播模型移出但仍需救治的病例。未受容量拒收限制的平均需求为
\[
D(t)=D_0(t)+\int_0^t h(t-u)
\frac{\beta c(u)S(u)I(u)}{N}\,du.
\]
利用非负性及总流入上界，有
\[
\begin{aligned}
D(t)
&\le D_{\mathrm{init}}+
\frac{\beta c_0S_0}{N}\eta\int_0^t h(a)\,da\\
&\le D_{\mathrm{init}}+pL\frac{\beta c_0S_0}{N}\eta,
\qquad t\ge0.
\end{aligned}
\]
设 $C>D_{\mathrm{init}}$ 为全过程固定的本病可用容量，已扣除其他病种占用和安全预留，本病初始病例的未来负担由 $D_{\mathrm{init}}$ 单独计入。于是
\[
\eta\le\frac{C-D_{\mathrm{init}}}{pL\,\beta c_0S_0/N}
\]
是平均需求满足 $D(t)\le C$ 的充分条件。该界对允许的干预分配共同成立，可能偏保守；它针对上述平均需求关系，不给出随机超载概率。按此条件选取阈值后，仍须检验正文中的触发条件和干预能力条件。

